% =========================================================
% PAGE 8: MACROECONOMIC OVERVIEW (ARABIC)
% =========================================================

% 1. Geometry (Sidebar on Right)
\GeomMacroMain


% ---------------------------------------------------------
% TITLE & INTRO
% ---------------------------------------------------------
\noindent
{\FontMacroTitle\textbf{\color{ecestitle}
\ECESContent{title_macro_overview_ar}{نظرة عامة على الاقتصاد الكلي}
}}

\vspace{1em}

\noindent
{\FontSubSection\textbf{
    \ECESContent{macro_global_intro_ar}{[أدخل: العنوان الرئيسي للمحور العالمي — ملخص في جملة واحدة لأبرز التطورات الاقتصادية العالمية.]}
    }}

\vspace{0.8em}

% ---------------------------------------------------------
% BODY TEXT
% ---------------------------------------------------------
\noindent
{\FontBody
\ECESContent{macro_global_p1_ar}{[أدخل: الفقرة الأولى — توقعات النمو العالمي ومعدلاته للاقتصادات المتقدمة والناشئة مع ذكر الأسباب والمصادر.]}

\vspace{0.8em}

\ECESContent{macro_global_p2_ar}{[أدخل: الفقرة الثانية — مسار التضخم العالمي وقرارات البنوك المركزية وأسعار الطاقة.]}

\vspace{0.8em}

\ECESContent{macro_global_p3_ar}{[أدخل: الفقرة الثالثة — مؤشر مديري المشتريات (PMI) ومستويات الثقة والطلب على العمالة.]}

\vspace{0.8em}

\textbf{\ECESContent{macro_local_p1_ar}{[أدخل: العنوان الرئيسي للمحور المحلي — ملخص في جملة واحدة لأبرز التطورات الاقتصادية المحلية.]}}

\vspace{0.8em}

\ECESContent{macro_local_p2_ar}{[أدخل: مراجعة صندوق النقد الدولي / البنك الدولي — التقييم والتوقعات الاقتصادية لمصر.]}

\vspace{0.8em}

\ECESContent{macro_local_p3_ar}{[أدخل: أولويات الإصلاحات الهيكلية الموصى بها لتعزيز النمو والاستقرار.]}
}

% ---------------------------------------------------------
% FOOTNOTES SECTION
% ---------------------------------------------------------
\vfill

\noindent
\rule{5cm}{0.4pt}

\vspace{0.3em}

\noindent
{\FontNote
\ECESContent{macro_footnotes_pg8_ar}{[أدخل: الحواشي والمراجع المتعلقة بمحتوى هذه الصفحة.]}
}

\clearpage

% =========================================================
% PAGE 9: MACROECONOMIC INDICATORS (ARABIC)
% =========================================================

% 1. Geometry (Sidebar on Right)
\GeomMacroMain


% ---------------------------------------------------------
% INTRO TEXT
% ---------------------------------------------------------
\noindent
{\FontBody
\ECESContent{macro_indicators_intro_ar}{[أدخل: جملة تمهيدية للجزء الخاص بمؤشرات الاقتصاد الكلي وأحدث البيانات المتاحة.]}
}

\vspace{0.8em}

% ---------------------------------------------------------
% SECTION I: GDP GROWTH
% ---------------------------------------------------------
\noindent
{\FontSubSection\textbf{\ECESContent{macro_heading_gdp_ar}{أولا: معدل نمو الناتج المحلي الإجمالي}}}

\noindent
{\FontBody
\ECESContent{macro_gdp_growth_ar}{[أدخل: معدل نمو الناتج المحلي الإجمالي للربع محل الدراسة مقارنة بالربع السابق والمناظر مع القطاعات المحركة للنمو والمثبطة له.]}
}

% ---------------------------------------------------------
% FIGURE 1: ch6.png
% ---------------------------------------------------------
\vspace{0.5em}
\ECESFigCaption{\ECESContent{macro_chart1_caption_ar}{الشكل 1: النمو في الناتج المحلي الإجمالي الحقيقي}}

\noindent
\includegraphics[width=\linewidth, keepaspectratio]{ch6.png}

\ECESSourceNote{\ECESContent{macro_chart1_source_ar}{المصدر: وزارة التخطيط والتنمية الاقتصادية والتعاون الدولي.}}

\vspace{1.5em}

% ---------------------------------------------------------
% SECTION II: INFLATION
% ---------------------------------------------------------
\noindent
{\FontSubSection\textbf{\ECESContent{macro_heading_inflation_ar}{ثانيا: معدل التضخم}}}

\noindent
{\FontBody
\ECESContent{macro_inflation_p1_ar}{[أدخل: معدل التضخم خلال الربع محل الدراسة مقارنة بالربع السابق مع الأسباب والمصادر.]}

\vspace{0.8em}

\ECESContent{macro_inflation_p2_ar}{[أدخل: قرار لجنة السياسة النقدية للبنك المركزي — أسعار الفائدة وأسباب القرار.]}
}

% ---------------------------------------------------------
% FOOTNOTES
% ---------------------------------------------------------
\vfill

\noindent
\rule{5cm}{0.4pt}

\vspace{0.3em}

\noindent
{\FontNote
\ECESContent{macro_footnotes_pg9_ar}{[أدخل: الحواشي والمراجع المتعلقة بمحتوى هذه الصفحة.]}
}

\clearpage

% =========================================================
% PAGE 10: INFLATION CHART & FOREIGN TRANSACTIONS
% =========================================================

% 1. Geometry (Sidebar on Right)
\GeomMacroMain

% ---------------------------------------------------------
% FIGURE 2: ch7.png
% ---------------------------------------------------------
\ECESFigCaption{\ECESContent{macro_chart2_caption_ar}{الشكل 2: التضخم وأسعار الفائدة الأساسية}}

\noindent
\includegraphics[width=\linewidth, keepaspectratio]{ch7.png}

{\FontNote\textit{\ECESContent{macro_chart2_source_ar}{المصادر: البنك المركزي المصري؛ والجهاز المركزي للتعبئة العامة والإحصاء.}}}

\vspace{1.5em}

% ---------------------------------------------------------
% SECTION III: FOREIGN TRANSACTIONS
% ---------------------------------------------------------
\noindent
{\FontSubSection\textbf{\ECESContent{macro_heading_foreign_ar}{ثالثا: المعاملات الخارجية}}}

\vspace{0.5em}

\noindent
{\FontBody
\textbf{\ECESContent{macro_label_bop_ar}{ميزان المدفوعات:}} \ECESContent{macro_bop_p1_ar}{[أدخل: تطور ميزان المدفوعات خلال الفترة — فائض أم عجز وأسبابه.]}

\vspace{0.6em}

\ECESContent{macro_bop_p2_ar}{[أدخل: تطور حساب المعاملات الجارية — التحويلات والسياحة والصادرات والواردات.]}

\vspace{0.6em}

\ECESContent{macro_bop_p3_ar}{[أدخل: تفاصيل إضافية لميزان المدفوعات — إيرادات قناة السويس والواردات البترولية.]}

\vspace{0.6em}

\textbf{\ECESContent{macro_label_nfa_ar}{صافي الأصول الأجنبية:}} \ECESContent{macro_nfa_p1_ar}{[أدخل: تطور صافي الأصول الأجنبية للقطاع المصرفي خلال الفترة.]}

\vspace{0.6em}

\textbf{\ECESContent{macro_label_reserves_ar}{صافي الاحتياطيات الدولية:}} \ECESContent{macro_reserves_p1_ar}{[أدخل: قيمة صافي الاحتياطيات الدولية في نهاية الربع ومقارنتها بالربع السابق والمناظر.]}

\vspace{0.6em}

\ECESContent{macro_exchange_rate_p1_ar}{[أدخل: قيمة الجنيه أمام الدولار خلال الربع (\CurrentQuarterText) ومقارنتها بالربع السابق والمناظر (الشكل 3).]}
}

% ---------------------------------------------------------
% FOOTNOTES
% ---------------------------------------------------------
\vfill

\noindent
\rule{5cm}{0.4pt}

\vspace{0.3em}

\noindent
{\FontNote
\ECESContent{macro_footnotes_pg10_ar}{[أدخل: الحواشي والمراجع المتعلقة بمحتوى هذه الصفحة.]}
}

\clearpage

% =========================================================
% PAGE 11: PUBLIC FINANCE (ARABIC)
% =========================================================

% 1. Geometry (Sidebar on Right)
\GeomMacroDeep


% ---------------------------------------------------------
% FIGURE 3: ch8.png
% ---------------------------------------------------------
\ECESFigCaption{\ECESContent{macro_chart3_caption_ar}{الشكل 3: صافي الاحتياطيات الدولية وسعر الصرف}}

\noindent
\includegraphics[width=\linewidth, keepaspectratio]{ch8.png}

{\FontNote\textit{\ECESContent{macro_chart3_source_ar}{المصادر: البنك المركزي المصري؛ ووزارة المالية.}}}

\vspace{2em}

% ---------------------------------------------------------
% SECTION IV: PUBLIC FINANCE
% ---------------------------------------------------------
\noindent
{\FontSubSection\textbf{\ECESContent{macro_heading_finance_ar}{رابعا: المالية العامة}}}

\vspace{0.5em}

\noindent
{\FontBody
\ECESContent{macro_public_finance_intro_ar}{[أدخل: العجز الكلي في الموازنة العامة كنسبة من الناتج المحلي الإجمالي ومقارنته بالعام السابق.]}
}

\begin{itemize} \itemsep0.6em
    \item \FontBody\ECESContent{macro_public_finance_b1_ar}{[أدخل: تطور الإيرادات العامة — الإيرادات الضريبية وغير الضريبية.]}

    \item \FontBody\ECESContent{macro_public_finance_b2_ar}{[أدخل: تطور المصروفات العامة — الإنفاق الاجتماعي وفوائد الدين.]}

    \item \FontBody\ECESContent{macro_public_finance_b3_ar}{[أدخل: نسبة الدين إلى الناتج المحلي الإجمالي وتطورها.]}
\end{itemize}

% ---------------------------------------------------------
% FOOTNOTES
% ---------------------------------------------------------
\vfill

\noindent
\rule{5cm}{0.4pt}

\vspace{0.3em}

\noindent
{\FontNote
\ECESContent{macro_footnotes_pg11_ar}{[أدخل: الحواشي والمراجع المتعلقة بمحتوى هذه الصفحة.]}
}

\restoregeometry
\clearpage
